\documentclass[UTF8,12pt]{ctexart}
\usepackage[a4paper,margin=24mm]{geometry}
\usepackage{amsmath,amssymb,bm,graphicx,adjustbox}
\newcommand{\fitmath}[1]{\begin{adjustbox}{max width=\linewidth}$\displaystyle #1$\end{adjustbox}}
\setlength{\parindent}{0pt}
\setlength{\parskip}{.7em}
\sloppy
\allowdisplaybreaks
\begin{document}
\section*{1987 年考研数学三 · 参考答案}
\subsection*{1987年真题参考答案}

\subsection*{（试卷IV）}

\subsection*{一、判断题}

（1）×。 （2）√。 （3）×。 （4）√。 （5）√。


\subsection*{二、选择题}

（1）A。 （2）C。 （3）C。 （4）A。 （5）C。


\subsection*{三、计算题}

（1）\(\displaystyle e\)。 （2）\(\displaystyle \dfrac{\displaystyle 2}{\displaystyle x\sqrt{1+x^2}}\)。 （3）\(\displaystyle dz=\dfrac{\displaystyle -y\,dx+x\,dy}{\displaystyle x^2+y^2}\)。


（4）\(\displaystyle (\sqrt{2x-1}-1)e^{\sqrt{2x-1}}+C\)，其中 \(\displaystyle C\) 为任意常数。


四、（1）当 \(\displaystyle t=\dfrac{\displaystyle \pi}{\displaystyle 4}\) 时，\(\displaystyle S_1+S_2\) 最小。（2）当 \(\displaystyle t=0\) 时，\(\displaystyle S_1+S_2\) 最大。


五、\(\displaystyle f(x)=\displaystyle \sum_{n=0}^{\infty}\left(1-\dfrac{\displaystyle 1}{\displaystyle 2^{n+1}}\right)x^n\)，收敛区间为 \(\displaystyle (-1,\allowbreak 1)\)。


六、\(\displaystyle \dfrac{\displaystyle e}{\displaystyle 2}-1\)。


七、\(\displaystyle x=e^{-p^3}\)。


八、\(\displaystyle (x_1,\allowbreak x_2,\allowbreak x_3,\allowbreak x_4)^{\mathrm T}=k(-1,\allowbreak 2,\allowbreak 1,\allowbreak 0)^{\mathrm T}+(3,\allowbreak -8,\allowbreak 0,\allowbreak 6)^{\mathrm T}\)，其中 \(\displaystyle k\) 为任意常数。


九、


\[\fitmath{\displaystyle B=\begin{pmatrix}3&-8&-6\\2&-9&-6\\-2&12&9\end{pmatrix}.}\]


十、特征值 \(\displaystyle \lambda=1\)，特征向量为 \(\displaystyle k(0,\allowbreak 2,\allowbreak 1)^{\mathrm T}\)，其中 \(\displaystyle k\) 为任意非零常数。


十一、（1）


\[\fitmath{\displaystyle F(x)=\begin{cases}0,&x<1,\\0.2,&1\le x<2,\\0.5,&2\le x<3,\\1,&x\ge3.\end{cases}}\]


（2）\(\displaystyle \dfrac{\displaystyle \sqrt{2\pi}}{\displaystyle 2a}\)。


十二、（1）\(\displaystyle \dfrac{\displaystyle 2}{\displaystyle 5}\)。（2）\(\displaystyle 0.48557\)。


\subsection*{（试卷V）}

\subsection*{一、判断题}

（1）【同试卷IV 第一、（1）题】 （2）【同试卷IV 第一、（2）题】 （3）×。 （4）×。 （5）【同试卷IV 第一、（5）题】


\subsection*{二、选择题}

（1）【同试卷IV 第二、（1）题】 （2）【同试卷IV 第二、（2）题】 （3）【同试卷IV 第二、（3）题】 （4）【同试卷IV 第二、（4）题】 （5）C。


\subsection*{（试卷V）参考答案续}

三、（1）\(\displaystyle 1\)。 （2）【同试卷IV 第三、（2）题】 （3）【同试卷IV 第三、（3）题】 （4）\(\displaystyle 1\)。


（5）\(\displaystyle \dfrac{\displaystyle 1}{\displaystyle 4}\arctan\dfrac{\displaystyle x^2+1}{\displaystyle 2}+C\)，其中 \(\displaystyle C\) 为任意常数。


四、（1）当 \(\displaystyle t=\dfrac{\displaystyle 1}{\displaystyle 2}\) 时，\(\displaystyle S_1+S_2\) 最小。（2）当 \(\displaystyle t=1\) 时，\(\displaystyle S_1+S_2\) 最大。


五、【同试卷IV 第六题】


六、（1）边际成本为 \(\displaystyle \dfrac{\displaystyle dC(x)}{\displaystyle dx}=3+x\)。（2）边际收益为 \(\displaystyle \dfrac{\displaystyle dR(x)}{\displaystyle dx}=\dfrac{\displaystyle 50}{\displaystyle \sqrt x}\)。


（3）边际利润为 \(\displaystyle \dfrac{\displaystyle d(R-C)}{\displaystyle dx}=\dfrac{\displaystyle 50}{\displaystyle \sqrt x}-3-x\)。（4）收益对价格的弹性为 \(\displaystyle -1\)。


七、【同试卷IV 第八题】 八、【同试卷IV 第九题】 九、【同试卷IV 第十题】


十、


\[\fitmath{\displaystyle F(x)=\begin{cases}0,&x<1,\\0.2,&1\le x<2,\\0.5,&2\le x<3,\\1,&x\ge3,\end{cases}\quad E(X)=2.3,\quad D(X)=0.61.}\]


十一、【同试卷IV 第十二题】

\end{document}
